% ID: SEL-D0031-OA-E02
% ESTADO: validada. S032, cotejada con examen y criterios.
\textbf{a)} En $1$, el valor del primer tramo es $b$ y el límite
derecho es $a-2$. La continuidad exige $b=a-2$. Las derivadas
de los tramos son $b/(2-x)^2$ y $2ax-3$. Para que coincidan
las pendientes en $1$ hace falta $b=2a-3$. Al igualar resulta
$\boxed{a=1,\ b=-1}$, que cumple las dos condiciones.

\textbf{b)} Con $a=1$ y $b=2$, el primer tramo tiene derivada
positiva y crece en $(-\infty,1]$. En el segundo, $2x-3$ es
negativa entre $1$ y $3/2$ y positiva después. Además, en
$1$ la función salta de $2$ a un límite derecho de $-1$.
Así, decrece en $[1,3/2]$ incluyendo el salto inicial y
crece en $[3/2,\infty)$.

Cuando $x\to-\infty$, $2/(2-x)\to0$. Por tanto $y=0$ es
asíntota horizontal por la izquierda. El denominador del primer
tramo nunca se anula en su dominio. En $1$ los límites son finitos,
y el crecimiento cuadrático por la derecha descarta una asíntota
recta allí. No hay otras asíntotas.
